\scanpage{47}
% archive foliation: 46
\[
\alpha_*(C_\bullet^i)\simeq DP(C_\bullet^i)\simeq\bigwedge^{i+1}\Phi_*
\]
\[
(\text{d'où}\qquad \alpha^*(C_i^\bullet)\simeq DP(C_i^\bullet)\simeq\bigwedge^{i+1}\unsure{\Phi_*|\Phi^*}\quad)
\]
\deleted{\unsure{L'obj. déf.} de $\alpha(C_\bullet^i)$ est \unsure{l'image} de $\unsure{T}_{*\wedge}$
par le foncteur $ND:K_*\longrightarrow K_\bullet$ \unsure{donné} \unsure{can.}
de $(\bigwedge^{\geqslant1})\unsure{\Phi}_*$.}
\deleted{\unsure{Par suite}}
\deleted{\uline{Cor. 1} Le foncteur $DP:Ab_{k\bullet}\longrightarrow\unsure{Ab}$}

Appliquant la théorie de Dold-Puppe à $K$
\deleted{(qui est bien additive et stable par facteurs
directs, au moins si $k$ est principal, p.ex.
$k=\mathbf Z$)} on trouve dans $K$ un complexe
\[
\deleted{\unsure{C}}\qquad C^*=DP(C^\bullet)
\]
(\unsure{qu'on} \deleted{\unsure{pourra prendre}} \unsure{a donc défini} dans \deleted{$\unsure{K}$} $K(k)$ \unsure{quelque}
\unsure{part soit} $k$), d'où pour les $B$ $k$-additives
\diagram{\begin{tikzcd}[row sep=large,column sep=large]
\underline{\operatorname{Hom}}_{k\text{-lin}}(K,B)\arrow[d,"\simeq"',"\mathrm{rest}"]
\arrow[ddr,bend left=18,"F\longmapsto F(C^*)"]&\\
\underline{\operatorname{Hom}}_{k\text{-lin}}(\unsure{A},B)\arrow[d,"\simeq"',"F\longmapsto F(C^\bullet)"]&\\
B^\bullet\arrow[r,"DP",hook]&B^*
\end{tikzcd}}
donc un foncteur $F\longmapsto F(C^*)$ qui est pl. fid.
Si $B$ est karoubienne, $DP_B$ est une équivalence,
donc aussi $F\longmapsto F(C^*)$. Donc
$C^*$ apparait comme objet simplicial \uline{universel}
dans cat. $k$-additive karoubienne.

\uline{Prop 5} Si $B$ est $k$-additive karoubienne, alors
$F\longmapsto F(C^*)$ est une équivalence
\[
\underline{\operatorname{Hom}}_{k\text{-lin}}(K,B)\longrightarrow B^*.
\]
(On a un \unsure{hom|sens} canoniquement \inserted{\unsure{bif. de}}
$F\in\operatorname{Ob}\underline{\operatorname{Hom}}_{k\text{-lin}}(K,B)$ et $x\in\operatorname{Ob}K$.
\margin{\unsure{i}}

